\appendix
\section{Author-Reported Hardware Results}
\label{sec:reported-results}
The authors supplied the following aggregate report and subsequently confirmed that its values are actually measured results, correcting its initial ``reference result'' label. We preserve that provenance distinction: these are author-reported measurements, separate from the locally executed and raw-log-verified experiments in the main text. Run-level CSV files, raw hardware logs, hashes and configuration manifests were not included with the report. Consequently, we can check arithmetic but cannot independently recompute these aggregates, verify the execution claims or derive statistical uncertainty. The report is included without upgrading these limitations into verified evidence.

\paragraph{Reported configuration and execution.}
The report specifies Linux, eBPF plus userspace telemetry sampled every 5\,s, Llama-3.1-8B-Instruct, four workload categories (CPU, memory, I/O, mixed), seven methods and 20 independent runs per workload/method. A complete matrix would contain 560 runs per server; the corresponding run inventory was not supplied. The authors report separate training, validation, calibration and test families, excluding test traces and graph edges from training and retrieval. Exact hardware, kernel release, checkpoint digest, code commit, load settings, budgets and split manifests are needed to make this reported setup reproducible. The supplied description does not identify the executable version used for deployed actuation; the artifact in this revision remains dry-run only.

\begin{table}[t]
\centering\small
\caption{Author-reported measured aggregate values. Relative changes are computed from the supplied aggregates; the aggregation rule and run-level uncertainty are unavailable. Anomaly reduction is 1.7 percentage points.}
\label{tab:reference-performance}
\begin{tabular}{lrrr}
\toprule
Metric & Default & Full & Change \\
\midrule
P50 (ms) & 8.5 & 8.1 & $-4.7\%$ \\
P95 (ms) & 20.0 & 17.2 & $-14.0\%$ \\
P99 (ms) & 35.0 & 30.2 & $-13.7\%$ \\
Throughput (req/s) & 10,000 & 10,280 & $+2.8\%$ \\
Anomaly (\%) & 3.2 & 1.5 & $-53.1\%$ \\
\bottomrule
\end{tabular}
\end{table}
\begin{table}[t]
\centering\small
\caption{Author-reported measured method comparison. Reduction is relative to the reported 20.0 ms Default value. Run-level uncertainty and matched-budget metadata were not supplied.}
\label{tab:reference-methods}
\begin{tabular}{lrr}
\toprule
Method & P95 (ms) & Reduction \\
\midrule
Default & 20.0 & --- \\
Rules + Graph & 19.0 & 5.0\% \\
BO + Graph & 18.2 & 9.0\% \\
RL + Graph & 18.0 & 10.0\% \\
LLM-only & 18.5 & 7.5\% \\
Graph + LLM & 17.4 & 13.0\% \\
Full SemantOS & 17.2 & 14.0\% \\
\bottomrule
\end{tabular}
\end{table}
\paragraph{Method comparisons.}
The reported Rules + Graph to Graph + LLM reduction is $(19.0-17.4)/19.0=8.4\%$, and the LLM-only to Graph + LLM reduction is $(18.5-17.4)/18.5=5.9\%$. Against the reported BO + Graph and RL + Graph values, Graph + LLM is respectively 4.4\% and 3.3\% lower. These aggregate contrasts are consistent with a benefit from the combined method but do not by themselves isolate a causal LLM reasoning effect: matched information, budgets, policy differences and paired run-level uncertainty remain to be documented. The full-system 14\% difference must not be attributed entirely to the LLM.

\begin{table}[t]
\centering\small
\caption{Author-reported measured workload P95 values (ms). The supplied Graph label has not been mapped to a specific implementation. These values do not reproduce the aggregate table by unweighted averaging.}
\label{tab:reference-workloads}
\begin{tabular}{lrrrrr}
\toprule
Workload & Default & Graph & LLM & G+L & Full \\
\midrule
CPU & 14.2 & 13.9 & 13.7 & 13.2 & 13.1 \\
Memory & 25.4 & 23.6 & 24.0 & 21.9 & 21.8 \\
I/O & 31.8 & 29.9 & 30.7 & 28.6 & 28.5 \\
Mixed & 22.1 & 21.2 & 20.8 & 19.3 & 19.1 \\
\bottomrule
\end{tabular}
\end{table}
\paragraph{Unresolved aggregation.}
The workload table implies full-system reductions of 7.7\%, 14.2\%, 10.4\% and 13.6\%. Its unweighted means are 23.375\,ms (Default) and 20.625\,ms (Full), rather than the aggregate table's 20.0 and 17.2\,ms. The ratio of those means gives 11.8\%, while averaging workload-level reductions gives 11.5\%. Neither reproduces 14.0\%. An alternative weighting or request-level aggregation may explain the difference, but no rule was supplied; averaging P95 values is not equivalent to a pooled request P95. We therefore retain the two reported tables separately and do not synthesize one headline estimate.

\begin{table}[t]
\centering\small
\caption{Author-reported measured safety percentages. Counts, denominators and independent unsafe labels were not supplied; these rates alone cannot establish veto recall or rollback precision.}
\label{tab:reference-safety}
\begin{tabular}{lrrr}
\toprule
Method & Unsafe (\%) & Veto (\%) & Rollback (\%) \\
\midrule
Rules + Graph & 8.0 & 5.1 & 1.4 \\
BO + Graph & 12.5 & 8.8 & 2.0 \\
RL + Graph & 11.0 & 7.3 & 1.9 \\
LLM-only & 10.0 & 6.7 & 1.8 \\
Graph + LLM & 5.5 & 3.8 & 1.0 \\
Full SemantOS & 4.0 & 3.1 & 0.6 \\
\bottomrule
\end{tabular}
\end{table}
\paragraph{Reported control-loop trace.}
The supplied trace reports initial P95/P99 of 18.9/31.4\,ms, throughput 11,240\,req/s and CPU utilization 73.1\%. It records recommendations of swappiness $60\to40$ and dirty ratio $20\to15$ with uncertainty 0.21, passing validation, followed by canary/ramp/full P95 values of 17.4/16.9/16.8\,ms and final P99 27.6\,ms without rollback. The report also describes range rejection of swappiness $-20$, and rollback when canary P95 rises from 18.2 to 25.4\,ms against a 22.0\,ms SLO, followed by P95 18.3\,ms. These are author-reported execution excerpts; raw logs, setting readbacks and implementation commit are needed to verify the deployed path and restoration. They do not establish that the current dry-run artifact performs these writes.

Veto precedes deployment; rollback follows an applied change. Their counts and denominators must remain distinct, with independently labeled unsafe actions and unknown labels retained. Nonzero veto or rollback rates alone do not prove effective protection. The reported reduction in unsafe proposals between Graph + LLM and Full SemantOS also requires a description of changes to candidate generation or feedback, because enforcement alone need not change proposed actions.

\paragraph{Reported uncertainty and repetition summary.}
The report associates scores 0.12 and 0.28 with acceptance, 0.57 with canary/delay, and 0.81 with veto for progressively less familiar workloads. The frozen threshold, calibration sample and implemented delay state were not supplied, so this ordering is not evidence of calibration. Its separate repetition summary gives $20.1\pm0.9$ versus $17.4\pm0.7$\,ms, differing from the aggregate table, without defining $\pm$. Six run values are listed (17.12, 17.48, 16.91, 17.35, 17.61, 17.09); twenty complete run records are not present. No confidence interval or significance test is reconstructed from these partial summaries. Identical rounded observations alone would not invalidate measurements, and variation alone would not verify them.

\paragraph{Evidence integration.}
The report specifies provenance fields for run ID, UTC timestamp, workload family, server ID, kernel release, code commit, method, repetition, split, latency quantiles, throughput, anomaly/window counts, proposal/veto/rollback counts and raw-log SHA-256. These fields describe the expected records; actual records and hashes were not supplied. Retaining proposal acceptance counts, independent anomaly definitions and exposure windows is also necessary. Once available, raw logs should produce run records, aggregates and paper tables through a reproducible pipeline. Until then, the reported reductions are descriptive, author-supplied aggregates with unresolved aggregation and uncertainty, rather than independently reproduced comparative estimates. Chat citation placeholders in the supplied report are not resolvable bibliographic references and are not treated as sources.
